\documentclass[presentation,aspectratio=169]{beamer}
\usepackage{tikz}
\usetikzlibrary{positioning,calc}
\usetikzlibrary{shapes.geometric}
\usetikzlibrary{backgrounds}% only to show the bounding box
\usetikzlibrary{shapes,arrows}
\usepackage{appendixnumberbeamer}
\usepackage{amsmath}
\newif\ifwidescreen
\widescreentrue
\date{28th May 2018}
\usetheme{metropolis}

\metroset{background=light,progressbar=frametitle,numbering=counter}

\renewcommand{\vec}[1]{\ensuremath{\boldsymbol{#1}}}
\newcommand{\ddt}[1]{\frac{\partial #1}{\partial t}}
\newcommand{\zhat}{\hat{\vec{z}}}
\newcommand{\W}{\ensuremath{\mathbb{W}}}

\newcommand{\inner}[1]{\left\langle #1 \right \rangle}

\newcommand{\KSP}[2]{\ensuremath{\mathcal{K}\left(#1, \mathbb{#2}\right)}}
\newcommand{\ksp}[1]{\KSP{#1}{#1}}

\newcommand{\colourfiredrake}[1]{\colorbox{red!20}{#1}}
\newcommand{\colourpetsc}[1]{\colorbox{blue!20}{#1}}

\author{Lawrence Mitchell\inst{1,*} \and Rob Kirby\inst{2,\dagger}}
\institute{
\inst{1}Departments of Computing and Mathematics, Imperial College
London

\inst{*}\texttt{lawrence.mitchell@imperial.ac.uk}
\and
\inst{2}Department of Mathematics, Baylor University

\inst{\dagger}\texttt{robert\_kirby@baylor.edu}
}

\graphicspath{{./\jobname.figures/}{../pictures/}}
\usepackage[url=false,
            doi=true,
            isbn=false,
            style=authoryear,
            firstinits=true,
            uniquename=init,
            backend=biber]{biblatex}

\setbeamertemplate{bibliography item}{}
\renewcommand{\bibfont}{\fontsize{7}{7}\selectfont}
\addbibresource{../literature.bib}
\newcommand{\arxivlink}[2]{%
  \href{http://www.arxiv.org/abs/#1}%
  {{\small\texttt{arXiv:\,#1\,[#2]}}}%
}

\setlength{\bibitemsep}{1ex}
\setlength{\fboxsep}{1pt}

\renewbibmacro{in:}{}
\DeclareFieldFormat[article]{volume}{\textbf{#1}}
\DeclareFieldFormat{doi}{%
  doi\addcolon%
  {\ifhyperref{\href{http://dx.doi.org/#1}{\nolinkurl{#1}}}
    {\nolinkurl{#1}}}}
\AtEveryBibitem{%
\clearfield{pages}%
\clearfield{issue}%
\clearfield{number}%
}

\usepackage{minted}
\title{Solver composition across the PDE/linear algebra divide}

\titlegraphic{\hfill\includegraphics[height=1.1cm]{firedrake-word.pdf}}

\begin{document}
\bgroup
\setbeamertemplate{title graphic}{
  \vbox to 0pt {
    \vspace*{0.25em}
    \makebox[1.025\textwidth]{\inserttitlegraphic}%
  }%
  \nointerlineskip%
}
\setbeamertemplate{background}{%
  \raisebox{-\paperheight}[0pt][0pt]{%
    \makebox[\paperwidth][c]{%
      \includegraphics[width=\paperwidth]{bottom-swoosh}%
    }%
  }%
}
\begin{frame}
  \maketitle{}
\end{frame}
\egroup

\begin{frame}[fragile,t]
  \frametitle{UFL makes it easy to write complex PDEs}
  \begin{onlyenv}<1>
    \begin{columns}[T]
      \begin{column}{0.47\framewidth}
        \small
        \begin{block}{Rayleigh-B\'enard convection}
          \begin{equation*}
            \begin{split}
              -\Delta u + u\cdot\nabla u + \nabla p +
              \frac{\text{Ra}}{\text{Pr}} \hat{g}T &= 0 \\
              \nabla \cdot u &= 0 \\
              - \frac{1}{\text{Pr}} \Delta T + u\cdot \nabla T &= 0
            \end{split}
          \end{equation*}
          Newton
          \begin{equation*}
            \begin{bmatrix}
              F   & B^T & M_1 \\
              C   & 0   & 0   \\
              M_2 & 0 & K
            \end{bmatrix}
            \begin{bmatrix}
              \delta u \\
              \delta p \\
              \delta T
            \end{bmatrix} =
            \begin{bmatrix}
              f_1 \\
              f_2 \\
              f_3
            \end{bmatrix}
          \end{equation*}
        \end{block}
      \end{column}
      \begin{column}{0.52\framewidth}
\begin{minted}[fontsize=\tiny]{python}
from firedrake import *
mesh = Mesh(...)
V = VectorFunctionSpace(mesh, "CG", 2)
W = FunctionSpace(mesh, "CG", 1)
Q = FunctionSpace(mesh, "CG", 1)
Z = V * W * Q
Ra = Constant(200)
Pr = Constant(6.18)
upT = Function(Z)
u, p, T = split(upT)
v, q, S = TestFunctions(Z)
bcs = [...] # no-flow + temp gradient
nullspace = MixedVectorSpaceBasis(
   Z, [Z.sub(0), VectorSpaceBasis(constant=True),
       Z.sub(2)])
F = (inner(grad(u), grad(v))
     + inner(dot(grad(u), u), v)
     - inner(p, div(v))
     + (Ra/Pr)*inner(T*g, v)
     + inner(div(u), q)
     + inner(dot(grad(T), u), S)
     + (1/Pr) * inner(grad(T), grad(S)))*dx

solve(F == 0, upT, bcs=bcs, nullspace=nullspace)
\end{minted}
      \end{column}
    \end{columns}
  \end{onlyenv}
  \begin{onlyenv}<2>
    \begin{columns}[T]
      \begin{column}{0.47\textwidth}
        \begin{block}{Ohta--Kawasaki}
          \small
          \begin{align*}
            u_t - \Delta w + \sigma(u - m) &= 0\\
            w + \epsilon^2 \Delta u - u(u^2 - 1) &= 0
          \end{align*}
          Implicit timestepping + Newton
          \begin{equation*}
            \begin{bmatrix}
              (1 + \Delta t \theta \sigma)M  & \Delta t\theta K \\
              -\epsilon^2 K - M_E & M
            \end{bmatrix}
            \begin{bmatrix}
              \delta u \\
              \delta w
            \end{bmatrix} =
            \begin{bmatrix}
              f_1 \\
              f_2
            \end{bmatrix}
          \end{equation*}
        \end{block}
      \end{column}
      \begin{column}{0.52\textwidth}
\begin{minted}[fontsize=\tiny,escapeinside=||]{python}
from firedrake import *
mesh = Mesh(...)
V = FunctionSpace(mesh, "CG", 1)
Z = V*V
|$\epsilon$| = Constant(0.02)
|$\sigma$| = Constant(100)
dt = Constant(eps**2)
|$\theta$| = Constant(0.5)
v, q = TestFunctions(Z)
z = Function(Z)
z0 = Function(Z)
u, w = split(z)
u0, w0 = split(z0)
u|$_\theta$| = (1 - |$\theta$|)*u0 + |$\theta$|*u
w|$_\theta$| = (1 - |$\theta$|)*w0 + |$\theta$|*w
dfdu = u**3 - u
F = ((u - u0)*v
     + dt*dot(grad(w|$_\theta$|), grad(v))
     + dt*|$\sigma$|*(u|$_\theta$| - m)*v
     + w*q - dfdu*q
     - |$\epsilon$|**2*dot(grad(u), grad(q)))*dx
while t < ...:
    z0.assign(z)
    solve(F == 0, z)
\end{minted}
      \end{column}
    \end{columns}
  \end{onlyenv}
\end{frame}

\begin{frame}[t]
  \frametitle{What about the solvers?}
  \begin{block}{Direct}
    \begin{itemize}
    \item Great for small problems;
    \item wait forever for large problems.
    \end{itemize}
  \end{block}
  \begin{block}{Iterative}
    \begin{itemize}
    \item Need good preconditioners;
    \item For many problems, algebraic manipulation of the operator is
      insufficient.
    \end{itemize}
  \end{block}
\end{frame}

\begin{frame}
  \frametitle{A wishlist}
  \begin{itemize}
  \item Access to standard block preconditioners
  \item Easy specification of auxiliary operators
  \item ``Simple'' configuration
  \item Arbitrary nesting: smoothers inside splits, multigrid,
    etc\ldots
  \item \emph{Easy} to extend
  \end{itemize}

  \begin{block}{Claim}
    \vspace{0.2ex}
    All of these things are possible (straightforward?) if the solver
    library can call back to the PDE library to create operators.
  \end{block}
\end{frame}

\begin{frame}
  \frametitle{Idea}
  \begin{itemize}
  \item Endow discretised operators with PDE-level information:
    \begin{itemize}
    \item what bilinear form
    \item which function spaces
    \item boundary conditions
    \end{itemize}
  \item Enable standard fieldsplitting on these operators.
  \item Write custom preconditioners that can utilise the information
    in appropriately.
  \end{itemize}
  \begin{block}<2->{Extend PETSc with Firedrake-level PCs}
  \begin{itemize}
  \item PETSc already provides \emph{algebraic} composition of solvers. \nocite{Brown:2012}

  \item Firedrake can provide auxiliary operators

  \item We just need to combine these appropriately.
  \end{itemize}
  \end{block}
\end{frame}

% \begin{frame}
%   \frametitle{Block preconditioning}
%   \begin{itemize}
%   \item Coupled problems are (typically) not amenable to black
%     box preconditioning.
%   \item<2-> Solution: block preconditioning
%   \end{itemize}
%   \begin{columns}<2>[t]
%     \begin{column}{0.5\textwidth}
%       \begin{block}{Schur complement}
%         \begin{equation*}
%           T = \begin{bmatrix}
%             A & 0 \\
%             0 & C A^{-1} B^T
%           \end{bmatrix}^{-1}
%           \begin{bmatrix}
%             A & B^T \\
%             C & 0
%           \end{bmatrix}
%         \end{equation*}
%         has minimal polynomial $T(T - I)(T^2 - T - I) = 0$.
%         \nocite{Murphy:2000,Ipsen:2001,Benzi:2005}
%       \end{block}
%     \end{column}
%     \begin{column}{0.5\textwidth}
%       \begin{block}{Function space}
%         If $\mathcal{A} : W \rightarrow W^*$, then
%         \begin{equation*}
%           \langle u, v \rangle_W^{-1} \mathcal{A}
%         \end{equation*}
%         has mesh independent condition number.
%         \nocite{Malek:2014,Mardal:2011,Kirby:2010}
%       \end{block}
%     \end{column}
%   \end{columns}
% \end{frame}

\begin{frame}[fragile,t]
  \frametitle{Implementation: two parts}
  \begin{block}{A new matrix type}
    \vspace{0.2ex}
    A shell matrix that implements matrix-free actions, and contains
    the symbolic information about the bilinear form.
    \begin{columns}
      \begin{column}{0.5\textwidth}
      \begin{equation*}
        y \leftarrow A x
      \end{equation*}
    \end{column}
    \hspace{-0.1\textwidth}
    \begin{column}{0.7\textwidth}
\begin{minted}[fontsize=\scriptsize]{python}
A = assemble(a, mat_type="matfree")
\end{minted}
    \end{column}
    \end{columns}
    \begin{flushright}
      \footnotesize Could do this all with assembled matrices if desired.
    \end{flushright}
  \end{block}
  \begin{block}{Custom preconditioners}
    \vspace{0.2ex}
    These matrices do not have entries, we create preconditioners
    that inspect the UFL and do the appropriate thing.
    \begin{columns}
      \begin{column}{0.5\textwidth}
    \begin{equation*}
      y \leftarrow \tilde{A}^{-1} x
    \end{equation*}
  \end{column}
  \hspace{-0.1\textwidth}
  \begin{column}{0.7\textwidth}
\begin{minted}[fontsize=\scriptsize]{python}
solve(a == L, x,
      {"mat_type": "matfree",
       "pc_type": "python",
       "pc_python_type": "AssembledPC"})
\end{minted}
  \end{column}
  \end{columns}
  \end{block}
\end{frame}
\begin{frame}[fragile]
  \frametitle{This sounds like hard work}
  Fortunately, \texttt{petsc4py} makes it easy to write these PCs.
\begin{minted}[fontsize=\tiny,escapeinside=||,mathescape=true]{python3}
class MyPC(object):
    def setUp(self, pc):
        A, P = pc.getOperators()
        # A and P are shell matrices, carrying the symbolic
        # discretisation information.
        # So I have access to the mesh, function spaces, etc...
        # Can inspect options dictionary here
        # do whatever
    def apply(self, pc, r, e):
        # Compute approximation to error given current residual
        # $e \gets A^{-1} r$

solve(..., solver_parameters={"pc_type": "python",
                              "pc_python_type": "MyPC"})
\end{minted}
  PETSc manages all the splitting and nesting already.  So this does
  the right thing \emph{inside} multigrid, etc\ldots
\end{frame}

\begin{frame}[fragile,t]
  \frametitle{A simple example}
  Matrix-free actions with AMG on the assembled operator.
  \vspace{2\baselineskip}
  \begin{columns}
    \begin{column}[b]{0.5\textwidth}
\begin{minted}[fontsize=\tiny,escapeinside=||]{python}
a = u*v*dx + dot(grad(u), grad(v)*dx)
opts = {|\colourpetsc{”ksp\_type”: ”cg”}|,
        |\colourfiredrake{”mat\_type”: ”matfree”}|,
        |\colourfiredrake{”pc\_type”: ”python”}|,
        |\colourfiredrake{”pc\_python\_type”: ”AssembledPC”}|,
        |\colourpetsc{”assembled\_pc\_type”: ”hypre”}|}
solve(a == L, x, solver_parameters=opts)
\end{minted}
    \end{column}
    \ifwidescreen
    \begin{column}{0.5\textwidth}
    \else
    \hspace{-0.1\textwidth}
    \begin{column}{0.6\textwidth}
    \fi
      \begin{tikzpicture}[scale=0.7, every node/.style={transform shape},
        firedrake/.style={rounded rectangle,draw,black, fill=red!20, minimum size=6mm, inner sep=2pt},
        petsc/.style={rectangle,draw,black, minimum size=6mm,
          fill=blue!20, inner sep=2pt},
        arrow/.style={-stealth,line width=1}]
        \node[petsc] (ainv)
        {$A^{-1} \approx \mathcal{K}(A, A_p^{-1})$};
        \node[firedrake, below left=of ainv] (Ax) {$Ax$};
        \node[firedrake,below=of Ax] (ax)
        {\texttt{assemble(action(a, x))}};
        \node[petsc,below right=of ainv] (ap) {$A_p^{-1}$};
        \node[firedrake,below=of ap] (Ap)
        {$A_p \leftarrow \text{\texttt{assemble(a)}}$};
        \node[petsc,below=of Ap] (Apinv)
        {\texttt{hypre($A_p$)}};

        \draw[arrow] (ainv) -- (Ax); \draw[arrow] (Ax) -- (ax);
        \draw[arrow] (ainv) -- (ap); \draw[arrow] (ap) -- (Ap);
        \draw[arrow] (Ap) -- (Apinv);
      \end{tikzpicture}
    \end{column}
  \end{columns}
\end{frame}

\begin{frame}
  \frametitle{A more complicated example}
  \begin{onlyenv}<1> A preconditioner for the Ohta--Kawasaki equation
    \parencite{Farrell:2017}
    \begin{equation*}
      \begin{split}
        u_t - \Delta w + \sigma(u - m) &= 0\\
        w + \epsilon^2 \Delta u - u(u^2 - 1) &= 0
      \end{split}
    \end{equation*}
    Newton iteration at each timestep solves
    \begin{equation*}
      \begin{bmatrix}
        (1 + \Delta t \theta \sigma)M  & \Delta t\theta K \\
        -\epsilon^2 K - M_E & M
      \end{bmatrix}
      \begin{bmatrix}
        \delta u \\
        \delta w
      \end{bmatrix} =
      \begin{bmatrix}
        f_1 \\
        f_2
      \end{bmatrix}
    \end{equation*}
  \end{onlyenv}

  \begin{onlyenv}<2> Preconditioning strategy: {\scriptsize
      \begin{equation*}
        \begin{bmatrix}
          \left[(1 + \Delta t \theta \sigma)M\right]^{-1}  & 0 \\
          0 & S^{-1}
        \end{bmatrix}
        \begin{bmatrix}
          I & 0\\
          (\epsilon^2 K + M_E)\left[(1 + \Delta t \theta
            \sigma)M\right]^{-1} & I\\
        \end{bmatrix}
        \begin{bmatrix}
          (1 + \Delta t \theta \sigma)M  & \Delta t\theta K \\
          -\epsilon^2 K - M_E & M
        \end{bmatrix}.
      \end{equation*}
    } Where
    \begin{equation*}
      S = M + (\epsilon^2 K + M_E) \left[(1 + \Delta t\theta\sigma)M\right]^{-1} \Delta t \theta K
    \end{equation*}
    is inverted iteratively, preconditioned by
    \begin{equation*}
      S^{-1} \approx S_p^{-1} = \hat{S}^{-1}M\hat{S}^{-1}
    \end{equation*}
    with
    \begin{equation*}
      \hat{S} = M + \epsilon\sqrt{(\Delta t \theta)/(1+\Delta t \theta\sigma)} K.
    \end{equation*}
  \end{onlyenv}
\end{frame}

\begin{frame}[fragile,t]
  \frametitle{Implementation}
  \begin{columns}[T]
    \begin{column}{0.49\textwidth}
\begin{minted}[fontsize=\tiny,mathescape,escapeinside=||]{python}
class OKPC(PCBase):
    def setUp(self, pc):
        _, P = pc.getOperators()
        ctx = P.getPythonContext()
        # User information about $\Delta t$, $\theta$, etc...
        dt, |$\theta$|, |$\epsilon$|, |$\sigma$| = ctx.appctx["parameters"]
        V = ctx.a.arguments()[0].function_space()
        c = (dt * |$\theta$|)/(1 + dt * |$\theta$| * |$\sigma$|)
        w = TrialFunction(V)
        q = TestFunction(V)
        # $\hat{S} = \inner{q, w} + \epsilon\sqrt{c}\inner{\nabla q, \nabla w}$, $c = \frac{\Delta t \theta}{1 + \Delta t \theta \sigma}$
        op = assemble(inner(w, q)*dx + 
                      |$\epsilon$|*sqrt(c)*inner(grad(w), grad(q))*dx)
        self.ksp = KSP().create(comm=pc.comm)
        self.ksp.setOptionsPrefix(pc.getOptionsPrefix + "hats_")
        self.ksp.setOperators(op.petscmat, op.petscmat)
        self.ksp.setFromOptions()
        mass = assemble(w*q*dx)
        self.mass = mass.petscmat
\end{minted}
    \end{column}
    \begin{column}{0.49\textwidth}
\begin{minted}[fontsize=\tiny,escapeinside=||,mathescape=true]{python3}
    def apply(self, pc, x, y):
        t1, t2 = self.work
        # $t_1 \leftarrow \hat{S}^{-1}x$
        self.ksp.solve(x, t1)
        # $t_2 \leftarrow M t_1$
        self.mass.mult(t1, t2)
        # $y \leftarrow \hat{S}^{-1}t_2 = \hat{S}^{-1} M \hat{S}^{-1} x$
        self.ksp.solve(t2, y)
\end{minted}
    \end{column}
  \end{columns}
\end{frame}

\begin{frame}[fragile,t]
  \frametitle{Usage}
  \begin{columns}[T]
    \ifwidescreen
    \begin{column}{0.5\textwidth}
      \vspace{-1em}
    \else
    \hspace{-0.8cm}
    \begin{column}{0.35\textwidth}
    \fi
\begin{minted}[fontsize=\tiny,escapeinside=||]{py}
opts = {
  "snes_lag_preconditioner": -1,
  |\colourfiredrake{”mat\_type”: ”matfree”}|,
  |\colourpetsc{”ksp\_type”: ”gmres”}|,
  |\colourpetsc{”pc\_type”: ”fieldsplit”}|,
  |\colourpetsc{”pc\_fieldsplit\_type”: ”schur”}|,
  |\colourpetsc{”pc\_fieldsplit\_schur\_factorization\_type”: ”lower”}|,
  |\colourpetsc{”fieldsplit\_0”}|: {
     |\colourpetsc{”ksp\_type”: ”chebyshev”}|,
     "ksp_max_it": 10,
     |\colourfiredrake{”pc\_type”: ”python”}|,
     |\colourfiredrake{”pc\_python\_type”: ”firedrake.AssembledPC”}|,
     |\colourpetsc{”assembled\_pc\_type”: ”sor”}|},
  |\colourpetsc{”fieldsplit\_1”}|: {
     |\colourpetsc{”ksp\_type”: ”richardson”}|,
     "ksp_max_it": 2,
     |\colourfiredrake{”pc\_type”: ”python”}|,
     |\colourfiredrake{”pc\_python\_type”: ”OKPC”}|,
     |\colourpetsc{”hats\_ksp\_type”: ”richardson”}|,
     "hats_ksp_max_it": 5,
     |\colourpetsc{”hats\_pc\_type”: ”hypre”}|
 }

appctx = {"parameters": (dt, theta, eps, sigma)}
solve(L == 0, z, solver_parameters=opts, appctx=appctx)
\end{minted}
  \end{column}
  \ifwidescreen
  \begin{column}{0.55\textwidth}
  \else
  \begin{column}{0.65\textwidth}
  \fi
  \begin{tikzpicture}[scale=0.55, every node/.style={transform shape},
    firedrake/.style={rounded rectangle,draw,black, fill=red!20, minimum size=6mm, inner sep=2pt},
    petsc/.style={rectangle,draw,black, fill=blue!20, minimum size=6mm, inner sep=2pt},
    arrow/.style={-stealth,line width=0.8}]
    \node[petsc] (Jinv) {$J^{-1} \approx \mathcal{K}(J,
      J_p^{-1})$};
    \node[firedrake, below left=0.1cm and 0.3cm of Jinv.south west] (Jx) {$Jx$};
    \node[firedrake,below=0.5cm of Jx.south] (jx) {\texttt{assemble(action(J, x))}};
    \node[petsc,below right=0.5cm and -0.1cm of Jinv] (Jpinv) {$J_p^{-1} =%
      \left[%
        \begin{smallmatrix}
          A^{-1} & 0\\
          0 & S^{-1}
        \end{smallmatrix}%
      \right]
      \left[%
      \begin{smallmatrix}
        I & 0 \\
        -CA^{-1} & 0 \\
      \end{smallmatrix}%
      \right]
      $};
    \node[petsc, below left=0.8cm and -1cm of Jpinv] (Ainv)  {$A^{-1}\approx \mathcal{K}(A,A_p^{-1})$};
    \node[firedrake, below left=1cm and 1.5cm of Ainv] (Ax) {$Ax$};
    \node[firedrake, below=3cm of Ax] (ax) {\texttt{assemble(action(A, x))}};
    \node[firedrake, below=1.5cm of Ainv] (Ap)
    {$A_p \leftarrow \text{\raisebox{0.5pt}{\texttt{assemble(A\_p)}}}$};
    \node[petsc, below=of Ap] (Apinv) {\texttt{SOR($A_p$)}};

    \node[petsc, below=1.3cm of Jpinv.south east, anchor=north] (Sinv) {$S^{-1} \approx \mathcal{K}(S,
      S_p^{-1})$};
    \node[petsc, below left=2cm and 0.3cm of Sinv, anchor=north] (S) {$Sx = (M - C A^{-1} B)x$};
    \node[firedrake, below=of S] (Cx) {$Cx$};
    \node[firedrake, right=of Cx] (Mx) {$Mx$};
    \node[firedrake, left=of Cx] (Bx) {$Bx$};

    \node[firedrake, below=1.5cm of Cx] (actionx) {\texttt{assemble(action($\cdot$, x))}};

    \node[firedrake, below right=0.8cm and -0.2cm of Sinv, anchor=north] (Spinv) {$S_p^{-1} = \hat{S}^{-1} M
      \hat{S}^{-1}$};
    \node[petsc, below=3.25cm of Spinv] (Hatinv) {$\hat{S}^{-1}
      \approx \mathcal{K}(\hat{S}, \hat{S}_p^{-1})$};
    \node[firedrake, below=1.5cm of Hatinv] (Hat)
    {$\hat{S} \leftarrow \text{\texttt{assemble(hats)}}$};
    \node[petsc, below=1cm of Hat] (Hathypre) {\texttt{hypre($\hat{S}$)}};

    \draw[arrow] (Jinv.south west) -- (Jx);
    \draw[arrow] (Jx) -- (jx);
    \draw[arrow] (Jinv.south east) -- (Jpinv);

    \draw[arrow] (Jpinv.south west) -- (Ainv.north);
    \draw[arrow] (Ainv.south west) -- (Ax);
    \draw[arrow] (Ax) -- (ax);
    \draw[arrow] (Ainv) -- (Ap);
    \draw[arrow] (Ap) -- (Apinv);
    \draw[arrow] (Jpinv) -- (Sinv);
    \draw[arrow] (Sinv) -- (S);
    \draw[arrow] (S) -- (Ainv.south east);

    \draw[arrow] (S) -- (Bx);
    \draw[arrow] (S) -- (Cx);
    \draw[arrow] (S) -- (Mx);

    \draw[arrow] (Bx) -- (actionx);
    \draw[arrow] (Cx) -- (actionx);
    \draw[arrow] (Mx) -- (actionx);

    \draw[arrow] (Sinv) -- (Spinv);
    \draw[arrow] (Spinv) -- (Mx);
    \draw[arrow] (Spinv) -- (Hatinv);
    \draw[arrow] (Hatinv) -- (Hat);
    \draw[arrow] (Hat) -- (Hathypre);
  \end{tikzpicture}
    \end{column}
  \end{columns}
\end{frame}

\begin{frame}
  \frametitle{Rayleigh-B\'enard solver}
  {\small
  For each Newton step, solve
  \begin{equation*}
    \KSP{\begin{bmatrix}
        F & B^T & M_1\\
        C & 0 & 0 \\
        M_2 & 0 & K
      \end{bmatrix}}{J}
  \end{equation*}
  using a preconditioner from \textcite{Howle:2012}:
  \begin{equation*}
    \mathbb{J} =
    \begin{bmatrix}
      \KSP{\begin{bmatrix}
          F & B^T\\
          C & 0
        \end{bmatrix}}{N} & 0\\
      0 & I
    \end{bmatrix}
    \begin{bmatrix}
      I & 0 & -M_1\\
      0 & I & 0 \\
      0 & 0 & I
    \end{bmatrix}
    \begin{bmatrix}
      I & 0 & 0\\
      0 & I & 0\\
      0 & 0 &\ksp{K}
    \end{bmatrix}
  \end{equation*}
  with
  \begin{equation*}
    \mathbb{N} = \begin{bmatrix}
      F & 0 \\
      0 & \mathcal{K}(S_p, \KSP{L_p}{L}\,F_p \, \KSP{M_p}{M})
    \end{bmatrix}
    \begin{bmatrix}
      I & 0\\
      -C & I
    \end{bmatrix}
    \begin{bmatrix}
      \ksp{F} & 0 \\
      0 & I
    \end{bmatrix}
  \end{equation*}
  and
  \begin{equation*}
    S_p = -C \ksp{F} B^T.
  \end{equation*}

  \begin{flushright}
    {\scriptsize \textcite[\S B.4]{Kirby:2018} shows full solver configuration.}
  \end{flushright}
  }
\end{frame}
\begin{frame}
  \frametitle{Weak scaling}
  Limited by performance of algebraic solvers on subblocks.
  \begin{center}
    \scriptsize
  \begin{tabular}{c|c|c|c|c}
DoFs ($\times 10^{6}$) & MPI processes & Newton its & Krylov its & Time (s)\\
\hline
0.7405 & 24 & 3 & 16 & 31.7\\
2.973 & 96 & 3 & 17 & 43.9\\
11.66 & 384 & 3 & 17 & 56\\
45.54 & 1536 & 3 & 18 & 85.2\\
185.6 & 6144 & 3 & 19 & 167\\
  \end{tabular}

  \begin{uncoverenv}<2>
  \begin{tabular}{c|c|c|c|c}
    DoFs ($\times 10^{6}$) & \multicolumn{2}{|c|}{Navier-Stokes iterations} & \multicolumn{2}{|c}{Temperature iterations}\\
    & Total & per solve & Total & per solve\\
\hline
0.7405 & 329 & 20.6 & 107 & 6.7 \\
2.973  & 365 & 21.5 & 132 & 7.8 \\
11.66  & 373 & 21.9 & 137 & 8.1 \\
45.54  & 403 & 22.4 & 151 & 8.4 \\
185.6  & 463 & 24.4 & 174 & 9.2 \\
  \end{tabular}
\end{uncoverenv}
  \end{center}
\end{frame}

\begin{frame}
  \frametitle{Schwarz smoothers}
  This approach works well for block solvers, what else can I do?
  \begin{block}{Building blocks}
    \begin{itemize}
    \item Decomposition of mesh into patches\uncover<2->{: \alert{PETSc}}
    \item Operators on each patch\uncover<3->{: \alert{Firedrake}}
    \item Solvers for each patch\uncover<4->{: \alert{PETSc}}
    \item Boundary conditions\uncover<5->{: \alert{Homogeneous
          Dirichlet only for now}}
    \end{itemize}
  \end{block}
  \begin{uncoverenv}<6->
    \begin{block}{Idea}
      \begin{itemize}
      \item PETSc PC using \texttt{DMPlex} to provide patches
      \item Callback interface provides operator on each patch to
        \texttt{PCApply}
      \item Normal \texttt{KSP} on each patch to do the solve
      \end{itemize}
    \end{block}
  \end{uncoverenv}
\end{frame}

\begin{frame}[fragile]
  \frametitle{Patch definition}
  Patch described by set of entities on which dofs are free.
  \begin{block}{Builtin}
    \vspace{0.2ex}
    Specify patches by selecting:
    \begin{enumerate}
    \item Entities to iterate over (vertices, cells, \ldots);
    \item Adjacency relation that gathers ``free'' dofs.  Some builtin:
      \begin{itemize}
      \item[\texttt{star}] all dofs in star of entity
      \item[\texttt{vanka}] all dofs in closure of star of entity
      \end{itemize}
    \end{enumerate}
  \end{block}
  \begin{block}{User-defined}
    \vspace{0.2ex}
    Write short function to define patches ``by hand''.
  \end{block}
  \begin{block}{Implementation}
    \vspace{0.2ex}
    \url{github.com/wence-/ssc}, available in PETSc RSN.
  \end{block}
\end{frame}

\begin{frame}[fragile]
  \frametitle{Example: P2-P1 Stokes}
  \begin{onlyenv}<1>
    Monolithic multigrid with Vanka smoother on each level.
\begin{minted}[fontsize=\tiny]{python3}
solver_parameters = {
    "mat_type": "matfree",
    # Flex-gmres due to nonlinear PC (gmres as smoother)
    "ksp_type": "fgmres",
    "pc_type": "mg",
    "mg_levels": {
        "ksp_type": "gmres",
        "ksp_max_it": 2,
        "pc_type": "python",
        "pc_python_type": "ssc.PatchPC",
        "patch_pc_patch_construction_type": "vanka",
        "patch_pc_patch_construction_dim": 0,  # patches over vertices
        "patch_pc_patch_vanka_dim": 0,         # what entities are in the constraint space?
        "patch_pc_patch_exclude_subspace": 1,  # which subspace to exclude?
        "patch_pc_patch_sub_mat_type": "seqaij",
        "patch_sub_ksp_type": "preonly",
        "patch_sub_pc_type": "lu",
        "patch_sub_pc_factor_shift_type": "nonzero"
    },
    "mg_coarse_pc_type": "lu",
}
solve(F == 0, u, solver_parameters=solver_parameters)
\end{minted}
  \end{onlyenv}
  \begin{onlyenv}<2>
    Schur complement only requires change of options.
\begin{minted}[fontsize=\tiny,mathescape=true]{python3}
solver_parameters = {
    "mat_type": "matfree",
    "ksp_type": "gmres",
    "ksp_monitor": None,
    "pc_type": "fieldsplit",
    # Use $\text{diag}(A^{-1}, S^{-1})$ as PC
    "pc_fieldsplit_type": "schur",
    "pc_fieldsplit_schur_factorization_type": "diag",
    "fieldsplit_0": {
        # AMG on velocity block
        "ksp_type": "preonly",
        "pc_type": "python",
        "pc_python_type": "firedrake.AssembledPC",
        "assembled_pc_type": "hypre",
    },
    "fieldsplit_1": {
        # Inverse mass matrix to precondition $S$
        "ksp_type": "richardson",
        "pc_type": "firedrake.MassInvPC",
    }
}
solve(F == 0, u, solver_parameters=solver_parameters)
\end{minted}
  \end{onlyenv}
\end{frame}

\begin{frame}
  \frametitle{Conclusions}

  \begin{itemize}
  \item Composable solvers, using PDE library to easily
    develop complex block preconditioners.

  \item Model formulation decoupled from solver configuration.

  \item Automatically takes advantage of any improvements in both
    PETSc and Firedrake.

  \item Same approach works for Schwarz-like methods.
  \end{itemize}
  \begin{center}
    \url{www.firedrakeproject.org}
  \end{center}

  \begin{flushright}
    {\footnotesize
    \textcite{Kirby:2018} \arxivlink{1706.01346}{cs.MS}}
  \end{flushright}
  \begin{tikzpicture}[remember picture,overlay]
    \node[at=(current page.south west), anchor=south west] {\includegraphics[width=2.5cm]{epsrc-logo}};
  \end{tikzpicture}
\end{frame}


\appendix
\begin{frame}
  \frametitle{References}
  \printbibliography[heading=none]
\end{frame}

\end{document}

%%% Local Variables: 
%%% mode: latex
%%% TeX-master: t
%%% End: 
